% Options for packages loaded elsewhere
\PassOptionsToPackage{unicode}{hyperref}
\PassOptionsToPackage{hyphens}{url}
%
\documentclass[
  12pt,
]{article}
\usepackage{amsmath,amssymb}
\usepackage{iftex}
\ifPDFTeX
  \usepackage[T1]{fontenc}
  \usepackage[utf8]{inputenc}
  \usepackage{textcomp} % provide euro and other symbols
\else % if luatex or xetex
  \usepackage{unicode-math} % this also loads fontspec
  \defaultfontfeatures{Scale=MatchLowercase}
  \defaultfontfeatures[\rmfamily]{Ligatures=TeX,Scale=1}
\fi
\usepackage{lmodern}
\ifPDFTeX\else
  % xetex/luatex font selection
\fi
% Use upquote if available, for straight quotes in verbatim environments
\IfFileExists{upquote.sty}{\usepackage{upquote}}{}
\IfFileExists{microtype.sty}{% use microtype if available
  \usepackage[]{microtype}
  \UseMicrotypeSet[protrusion]{basicmath} % disable protrusion for tt fonts
}{}
\makeatletter
\@ifundefined{KOMAClassName}{% if non-KOMA class
  \IfFileExists{parskip.sty}{%
    \usepackage{parskip}
  }{% else
    \setlength{\parindent}{0pt}
    \setlength{\parskip}{6pt plus 2pt minus 1pt}}
}{% if KOMA class
  \KOMAoptions{parskip=half}}
\makeatother
\usepackage{xcolor}
\usepackage[margin=2.5cm]{geometry}
\usepackage{longtable,booktabs,array}
\usepackage{calc} % for calculating minipage widths
% Correct order of tables after \paragraph or \subparagraph
\usepackage{etoolbox}
\makeatletter
\patchcmd\longtable{\par}{\if@noskipsec\mbox{}\fi\par}{}{}
\makeatother
% Allow footnotes in longtable head/foot
\IfFileExists{footnotehyper.sty}{\usepackage{footnotehyper}}{\usepackage{footnote}}
\makesavenoteenv{longtable}
\usepackage{graphicx}
\makeatletter
\newsavebox\pandoc@box
\newcommand*\pandocbounded[1]{% scales image to fit in text height/width
  \sbox\pandoc@box{#1}%
  \Gscale@div\@tempa{\textheight}{\dimexpr\ht\pandoc@box+\dp\pandoc@box\relax}%
  \Gscale@div\@tempb{\linewidth}{\wd\pandoc@box}%
  \ifdim\@tempb\p@<\@tempa\p@\let\@tempa\@tempb\fi% select the smaller of both
  \ifdim\@tempa\p@<\p@\scalebox{\@tempa}{\usebox\pandoc@box}%
  \else\usebox{\pandoc@box}%
  \fi%
}
% Set default figure placement to htbp
\def\fps@figure{htbp}
\makeatother
\setlength{\emergencystretch}{3em} % prevent overfull lines
\providecommand{\tightlist}{%
  \setlength{\itemsep}{0pt}\setlength{\parskip}{0pt}}
\setcounter{secnumdepth}{-\maxdimen} % remove section numbering
\ifLuaTeX
\usepackage[bidi=basic]{babel}
\else
\usepackage[bidi=default]{babel}
\fi
\babelprovide[main,import]{brazilian}
% get rid of language-specific shorthands (see #6817):
\let\LanguageShortHands\languageshorthands
\def\languageshorthands#1{}
\usepackage{booktabs}
\usepackage{longtable}
\usepackage{array}
\usepackage{multirow}
\usepackage{wrapfig}
\usepackage{float}
\usepackage{colortbl}
\usepackage{pdflscape}
\usepackage{tabu}
\usepackage{threeparttable}
\usepackage{threeparttablex}
\usepackage[normalem]{ulem}
\usepackage{makecell}
\usepackage{xcolor}
\usepackage{bookmark}
\IfFileExists{xurl.sty}{\usepackage{xurl}}{} % add URL line breaks if available
\urlstyle{same}
\hypersetup{
  pdftitle={Trabalho 1 - ME613 : Análise de Regressão},
  pdflang={pt-BR},
  hidelinks,
  pdfcreator={LaTeX via pandoc}}

\title{Trabalho 1 - ME613 : Análise de Regressão}
\author{\textbf{Lucas Lima}: RA: 263267 \textbf{Geovani Ariel}: RA:
236124 \textbf{Daniela Cielo}: RA: 298825}
\date{2º semestre de 2025}

\begin{document}
\maketitle

{
\setcounter{tocdepth}{2}
\tableofcontents
}
\section{Introdução e
Motivação}\label{introduuxe7uxe3o-e-motivauxe7uxe3o}

A expectativa de vida é um indicador fundamental para avaliar o nível de
saúde e desenvolvimento de uma população, refletindo não apenas as
condições médicas, mas também fatores socioeconômicos, ambientais e de
políticas públicas. Ao sintetizar o tempo médio de vida que uma coorte
hipotética de indivíduos teria, caso permanecessem expostos às taxas de
mortalidade observadas em determinado período, esse indicador fornece
uma visão abrangente das condições de vida em diferentes países e
contextos históricos.

Nas últimas décadas, o aumento da expectativa de vida tem sido um dos
fenômenos demográficos mais marcantes em grande parte do mundo,
especialmente em países desenvolvidos. No entanto, essa evolução não é
homogênea: muitas regiões ainda enfrentam desafios importantes
relacionados à mortalidade infantil, doenças infecciosas, desnutrição,
desigualdade de renda e acesso limitado a serviços básicos de saúde e
educação. Essas disparidades revelam que ganhos agregados em expectativa
de vida podem ocultar heterogeneidades relevantes entre países, grupos
populacionais e períodos de tempo.

Do ponto de vista de políticas públicas, compreender os determinantes da
expectativa de vida é essencial para orientar decisões estratégicas,
como priorização de investimentos em saúde, educação e infraestrutura
social. Além disso, a análise conjunta de variáveis de mortalidade,
indicadores econômicos e fatores sociais permite investigar em que
medida os avanços em desenvolvimento humano se traduzem, de fato, em
melhores condições de sobrevivência e bem-estar da população.

Nesse contexto, este trabalho busca analisar a relação entre a
expectativa de vida e um conjunto de variáveis explicativas associadas a
fatores de mortalidade, imunização, condições econômicas e fatores
sociais, utilizando um conjunto de dados em painel com países observados
ao longo do tempo. A partir de técnicas de análise exploratória e
modelagem de regressão, pretende-se identificar quais fatores apresentam
maior associação com a expectativa de vida e discutir implicações para
políticas públicas e planejamento em saúde.

O estudo também se justifica pela relevância crescente de indicadores
compostos e modelos estatísticos que permitam capturar heterogeneidade
entre diferentes regiões do mundo. Em vez de analisar apenas um
indicador isolado, como o PIB per capita, o presente trabalho considera
um conjunto mais rico de variáveis, incluindo escolaridade média,
composição de renda de recursos, gastos em saúde e cobertura vacinal,
permitindo uma visão mais abrangente dos determinantes da longevidade.

O conjunto de dados utilizado neste trabalho foi construído a partir de
bases internacionais que reúnem informações demográficas, econômicas e
de saúde para diversos países e anos. Após tratamento e harmonização,
resultou em um painel com medidas anuais de expectativa de vida e
fatores associados. A estrutura em painel possibilita examinar tanto
diferenças entre países quanto a evolução temporal dentro de cada país,
permitindo investigar tendências globais e regionais na expectativa de
vida e nos fatores que mais contribuem para sua variação.

Entre as variáveis de interesse destacam-se:

\begin{itemize}
\tightlist
\item
  \textbf{Fatores de imunização:} Hepatite B, Poliomielite e Difteria;\\
\item
  \textbf{Fatores de mortalidade:} mortalidade infantil, mortalidade
  adulta, entre outros;\\
\item
  \textbf{Fatores econômicos:} PIB per capita, índices de
  desenvolvimento, gastos públicos;\\
\item
  \textbf{Fatores sociais:} escolaridade média, infraestrutura e
  condições de vida;
\end{itemize}

A escolha dessas categorias busca fornecer uma visão ampla sobre os
determinantes da expectativa de vida, combinando variáveis diretamente
relacionadas à mortalidade com indicadores que capturam o contexto
socioeconômico e de políticas públicas de cada país. Dessa forma, este
estudo tem como principal motivação identificar padrões e relações que
possam auxiliar formuladores de políticas e gestores de saúde na
priorização de políticas públicas e intervenções capazes de melhorar a
qualidade de vida da população.

\begin{longtable}[]{@{}lccc@{}}
\caption{Primeiras 5 observações - Colunas selecionadas}\tabularnewline
\toprule\noalign{}
Expectativa\_Vida & Mortalidade\_Adulto & Escolaridade & Populacao \\
\midrule\noalign{}
\endfirsthead
\toprule\noalign{}
Expectativa\_Vida & Mortalidade\_Adulto & Escolaridade & Populacao \\
\midrule\noalign{}
\endhead
\bottomrule\noalign{}
\endlastfoot
65.0 & 263 & 10.1 & 33736494 \\
59.9 & 271 & 10.0 & 327582 \\
59.9 & 268 & 9.9 & 31731688 \\
59.5 & 272 & 9.8 & 3696958 \\
59.2 & 275 & 9.5 & 2978599 \\
\end{longtable}

Para uma análise exploratória inicial, a Tabela 1 apresenta uma amostra
das primeiras cinco observações do banco de dados, contendo a variável
resposta e exemplos de variáveis explicativas de cada categoria
analisada. Especificamente, são mostradas:

\begin{itemize}
\tightlist
\item
  \textbf{Expectativa de Vida (Expectativa\_Vida):} variável resposta do
  estudo, representando o desfecho principal a ser explicado pelos
  demais fatores;\\
\item
  \textbf{Mortalidade Adulta (Mortalidade\_Adulto):} exemplo de fator de
  mortalidade, refletindo o risco de morte em idades adultas;\\
\item
  \textbf{Escolaridade (Escolaridade):} representando fatores sociais,
  relacionada ao nível médio de instrução da população;\\
\item
  \textbf{População (Populacao):} variável demográfica de contexto, que
  pode influenciar tanto a organização dos sistemas de saúde quanto o
  impacto relativo de políticas públicas.
\end{itemize}

A seleção da expectativa de vida como variável resposta fundamenta-se em
sua relevância teórica e prática, bem como na sua capacidade de
sintetizar múltiplos aspectos do desenvolvimento humano. As demais
variáveis ilustradas nesta tabela representam dimensões complementares
--- mortalidade, capital humano e estrutura demográfica --- que serão
exploradas em profundidade nas análises subsequentes.

\section{Análise Descritiva}\label{anuxe1lise-descritiva}

\begin{longtable}[t]{rrrr}
\caption{\label{tab:unnamed-chunk-4}Resumo geral da base de dados}\\
\toprule
n\_obs & n\_paises & ano\_min & ano\_max\\
\midrule
2938 & 193 & 2000 & 2015\\
\bottomrule
\end{longtable}

\begin{longtable}[t]{lrr}
\caption{\label{tab:unnamed-chunk-5}Quantidade e percentual de valores faltantes por variável}\\
\toprule
Variavel & N\_NA & Perc\_NA\\
\midrule
Populacao & 652 & 22.19\\
Hepatite\_B & 553 & 18.82\\
PIB & 448 & 15.25\\
Despesa\_Total & 226 & 7.69\\
Alcool & 194 & 6.60\\
\addlinespace
Composicao\_Renda\_Recursos & 167 & 5.68\\
Escolaridade & 163 & 5.55\\
IMC & 34 & 1.16\\
Magreza\_1\_19\_Anos & 34 & 1.16\\
Magreza\_5\_9\_Anos & 34 & 1.16\\
\addlinespace
Poliomielite & 19 & 0.65\\
Difteria & 19 & 0.65\\
Expectativa\_Vida & 10 & 0.34\\
Mortalidade\_Adulto & 10 & 0.34\\
Pais & 0 & 0.00\\
\addlinespace
Ano & 0 & 0.00\\
Status\_Pais & 0 & 0.00\\
Mortalidade\_Infantil & 0 & 0.00\\
Percentual\_Despesa & 0 & 0.00\\
Sarampo & 0 & 0.00\\
\addlinespace
Mortes\_Menores\_Cinco\_Anos & 0 & 0.00\\
HIV\_AIDS & 0 & 0.00\\
\bottomrule
\end{longtable}

\begin{longtable}[t]{rrrrrrrr}
\caption{\label{tab:unnamed-chunk-6}Resumo descritivo da variável Expectativa_Vida}\\
\toprule
n & media & sd & min & q1 & mediana & q3 & max\\
\midrule
2928 & 69.22 & 9.52 & 36.3 & 63.1 & 72.1 & 75.7 & 89\\
\bottomrule
\end{longtable}

\begin{longtable}[t]{lrrrrrrrr}
\caption{\label{tab:unnamed-chunk-7}Resumo da Expectativa_Vida por Status do país}\\
\toprule
Status\_Pais & n & media & sd & min & q1 & mediana & q3 & max\\
\midrule
Developed & 512 & 79.20 & 3.93 & 69.9 & 76.8 & 79.25 & 81.7 & 89\\
Developing & 2416 & 67.11 & 9.01 & 36.3 & 61.1 & 69.00 & 74.0 & 89\\
\bottomrule
\end{longtable}

\pandocbounded{\includegraphics[keepaspectratio]{TRABALHO_FINAL_ME613_files/figure-latex/unnamed-chunk-8-1.pdf}}

\pandocbounded{\includegraphics[keepaspectratio]{TRABALHO_FINAL_ME613_files/figure-latex/unnamed-chunk-9-1.pdf}}

\begin{verbatim}
## Warning: `aes_string()` was deprecated in ggplot2 3.0.0.
## i Please use tidy evaluation idioms with `aes()`.
## i See also `vignette("ggplot2-in-packages")` for more information.
## i The deprecated feature was likely used in the ggcorrplot package.
##   Please report the issue at <https://github.com/kassambara/ggcorrplot/issues>.
## This warning is displayed once every 8 hours.
## Call `lifecycle::last_lifecycle_warnings()` to see where this warning was
## generated.
\end{verbatim}

\pandocbounded{\includegraphics[keepaspectratio]{TRABALHO_FINAL_ME613_files/figure-latex/unnamed-chunk-10-1.pdf}}

A análise descritiva realizada evidencia que a Expectativa\_Vida
apresenta valores bastante heterogêneos entre países e ao longo do
tempo, com amplitude ampla entre os menores e maiores níveis observados.
Mesmo assim, a distribuição concentra-se em torno de valores
intermediários a altos, sugerindo que a maior parte das observações
corresponde a países em estágios de desenvolvimento em que a transição
epidemiológica já avançou em alguma medida.

Quando a variável resposta é comparada entre os diferentes níveis de
Status\_Pais, verifica-se uma diferença clara entre países classificados
como desenvolvidos e em desenvolvimento. De modo geral, países
desenvolvidos apresentam expectativa de vida média mais alta e menor
variabilidade, enquanto países em desenvolvimento exibem média mais
baixa e maior dispersão. Esse padrão é consistente com o esperado
teoricamente, refletindo desigualdades em termos de infraestrutura de
saúde, condições socioeconômicas e acesso a políticas públicas efetivas.

A análise temporal mostra uma tendência de aumento da expectativa de
vida ao longo dos anos no conjunto da amostra. Esse crescimento sugere
avanços globais em cuidados de saúde, imunização, controle de doenças
infecciosas e melhoria de condições de vida. No entanto, a velocidade
desse aumento pode variar entre os grupos de países, indicando que parte
das desigualdades em saúde persiste ao longo do tempo.

Os resultados da matriz de correlação reforçam a interpretação de que a
expectativa de vida é um indicador fortemente associado tanto a fatores
de mortalidade quanto a fatores sociais e econômicos. Entre as
associações negativas mais fortes destacam-se as correlações com
Mortalidade\_Adulto, Mortalidade\_Infantil, Magreza\_1\_19\_Anos e
Magreza\_5\_9\_Anos, além de variáveis relacionadas à carga de doenças,
como HIV\_AIDS. Essas relações sugerem que contextos com maior carga de
mortalidade e desnutrição tendem a apresentar expectativa de vida
substancialmente menor, o que é coerente com a compreensão
epidemiológica tradicional.

Por outro lado, a expectativa de vida apresenta correlações positivas
importantes com Escolaridade, Composicao\_Renda\_Recursos, PIB e, em
menor medida, com variáveis associadas à cobertura vacinal e à condição
nutricional (como IMC). Esses resultados indicam que maior
desenvolvimento socioeconômico e maior capital humano estão associados a
maior longevidade, possivelmente por meio de melhor acesso a serviços de
saúde, maior capacidade de prevenção de doenças e melhores condições
gerais de vida.

Do ponto de vista da modelagem, a EDA sugere que: • Existem diversos
preditores com correlação forte ou moderada com a Expectativa\_Vida, o
que é promissor para a construção de modelos de regressão com bom poder
explicativo; • Algumas variáveis explicativas são entre si fortemente
correlacionadas, o que levanta a hipótese de multicolinearidade e
reforça a necessidade de uma etapa cuidadosa de seleção de variáveis e
diagnóstico do modelo; • A presença de valores faltantes em variáveis
importantes (como PIB, Populacao ou alguns indicadores de saúde) deve
ser tratada com cuidado, seja por meio de exclusão de casos, imputação
ou construção de modelos alternativos, para não introduzir viés nas
estimativas.

Em síntese, a análise exploratória confirma a relevância da
Expectativa\_Vida como indicador-síntese de desenvolvimento e aponta um
conjunto coerente de variáveis candidatas a preditoras, abrangendo
dimensões de mortalidade, imunização, condição econômica e fatores
sociais. Esses achados fornecem uma base sólida para a etapa seguinte do
trabalho, na qual serão propostos modelos de regressão com o objetivo de
quantificar a contribuição relativa de cada fator e avaliar o ajuste
global do modelo à realidade observada.

\end{document}
